% Part: sets-functions-relations
% Chapter: relations
% Section: special-properties

\documentclass[../../../include/open-logic-section]{subfiles}

\begin{document}

\olfileid{sfr}{rel}{prp}
\olsection{Bijzondere eigenschappen van relaties}

\begin{intro}
Sommige soorten relaties komen zo vaak voor dat ze een eigen naam
hebben gekregen. Zo verhouden de elementen van hun respectieve
domeinen zich onder $\le$ en~$\subseteq$ op vergelijkbare wijze (neem
bijvoorbeeld $\Nat$ voor~$\le$ en $\Pow{A}$ voor~$\subseteq$). Om
precies aan te geven waarin deze relaties overeenkomen en verschillen,
delen we ze in naar enkele bijzondere eigenschappen die relaties
kunnen hebben. Enkele van die eigenschappen, of combinaties ervan,
leveren bijzonder belangrijke soorten relaties op: ordeningen en
equivalentierelaties.
\end{intro}

\begin{defn}[Reflexiviteit]
Een relatie $R \subseteq A^2$ is \emph{reflexief} dan en slechts dan
als voor elke $x \in A$ geldt dat $Rxx$.
\end{defn}

\begin{defn}[Transitiviteit]
Een relatie $R \subseteq A^2$ is \emph{transitief} dan en slechts dan
als uit $Rxy$ en $Ryz$ steeds $Rxz$ volgt.
\end{defn}

\begin{defn}[Symmetrie]
Een relatie~$R \subseteq A^2$ is \emph{symmetrisch} dan en slechts dan
als uit $Rxy$ steeds ook~$Ryx$ volgt.
\end{defn}

\begin{defn}[Antisymmetrie]
Een relatie~$R \subseteq A^2$ is \emph{an\-ti\-symmetrisch} dan en
slechts dan als uit $Rxy$ en $Ryx$ samen volgt dat $x=y$ (anders
gezegd: als $x\neq y$, dan geldt $\lnot Rxy$ of $\lnot Ryx$).
\end{defn}

\begin{explain}
Bij een symmetrische relatie gelden $Rxy$ en $Ryx$ altijd beide, of
geen van beide. Bij een antisymmetrische relatie kunnen $Rxy$ en $Ryx$
alleen beide gelden als $x = y$. Dit \emph{vereist} niet dat $Rxy$ en $Ryx$ gelden
wanneer $x = y$; het sluit dat slechts niet uit. Een
antisymmetrische relatie kan dus reflexief zijn, maar hoeft dat niet te
zijn. Antisymmetrisch zijn is bovendien iets anders dan alleen niet
symmetrisch zijn. Een relatie kan zelfs tegelijk symmetrisch en
antisymmetrisch zijn (de identiteitsrelatie is hiervan een voorbeeld).
\end{explain}

\begin{defn}[Samenhangendheid]
Een relatie $R \subseteq A^2$ is \emph{samenhangend} als voor alle
$x,y\in A$ geldt: als $x \neq y$, dan $Rxy$ of~$Ryx$.
\end{defn}

\begin{prob}
Geef voorbeelden van relaties die (a) reflexief en symmetrisch maar
niet transitief zijn, (b) reflexief en antisymmetrisch zijn, (c)
antisymmetrisch en transitief maar niet reflexief zijn, en (d)
reflexief, symmetrisch en transitief zijn. Gebruik geen relaties op
getallen of verzamelingen.
\end{prob} 
 
\begin{defn}[Irreflexiviteit]
Een relatie $R \subseteq A^2$ heet \emph{irreflexief} als voor alle
$x \in A$ geldt dat $Rxx$ niet geldt.
\end{defn}

\begin{defn}[Asymmetrie]
Een relatie $R \subseteq A^2$ heet \emph{asymmetrisch} als er geen paar
$x,y\in A$ is waarvoor zowel $Rxy$ als~$Ryx$ geldt.
\end{defn}

Als $A \neq \emptyset$, is geen enkele irreflexieve relatie op~$A$
reflexief en is elke asymmetrische relatie op~$A$ ook
antisymmetrisch. Er zijn echter relaties $R \subseteq A^2$ die noch
reflexief noch irreflexief zijn, en er zijn antisymmetrische relaties
die niet asymmetrisch zijn.
\end{document}